\documentclass{article}%
\usepackage[T1]{fontenc}%
\usepackage[utf8]{inputenc}%
\usepackage{lmodern}%
\usepackage{textcomp}%
\usepackage{lastpage}%
\usepackage{geometry}%
\geometry{margin=1in,headheight=14pt,headsep=25pt}%
\usepackage{hyperref}%
\usepackage{enumitem}%
\usepackage{fancyhdr}%
\usepackage{xcolor}%
%

            \hypersetup{
                colorlinks=true,
                linkcolor=blue,
                filecolor=magenta,
                urlcolor=blue,
            }
            
            \pagestyle{fancy}
            \fancyhf{}
            \rhead{Generated on \today}
            \lhead{Course Summary}
            \cfoot{\thepage}
            
            % Configure itemize settings globally
            \setlist[itemize]{nosep}  % Reduce vertical spacing between items
            \setlist[itemize]{leftmargin=*}  % Align with left margin
        %
\title{Linear Programming Course Summary}%
\author{Generated by Scribe.AI}%
\date{\today}%
%
\begin{document}%
\normalsize%
\maketitle%
\section{Key Terms}%
\label{sec:KeyTerms}%
\begin{itemize}[leftmargin=*]%
\subsection{linear programming fundamentals}%
\label{subsec:linearprogrammingfundamentals}%
\begin{itemize}[leftmargin=*]%
\item \href{https://www.math.purdue.edu/~yipn/421/2024-08-27-ExSimplex.pdf#page=1}{\textbf{\textcolor{blue}{Slack Variable}}}: \emph{A variable added to a less-than-or-equal-to inequality constraint to transform it into an equality constraint. It represents the difference between the left-hand side and right-hand side of the inequality.}%
\end{itemize}

%
\subsection{solution space analysis}%
\label{subsec:solutionspaceanalysis}%
\begin{itemize}[leftmargin=*]%
\item \href{https://www.math.purdue.edu/~yipn/421/2024-08-27-ExSimplex.pdf#page=4}{\textbf{\textcolor{blue}{Feasible Region}}}: \emph{The set of all points that satisfy all the constraints of a linear programming problem.}%
\item \href{https://www.math.purdue.edu/~yipn/421/2024-08-27-ExSimplex.pdf#page=5}{\textbf{\textcolor{blue}{Vertex}}}: \emph{A point where two or more constraint boundaries intersect in the feasible region of a linear programming problem. In higher dimensions, it is the intersection of multiple constraint hyperplanes.}%
\end{itemize}

%
\subsection{simplex method algorithms}%
\label{subsec:simplexmethodalgorithms}%
\begin{itemize}[leftmargin=*]%
\item \href{https://www.math.purdue.edu/~yipn/421/2024-08-27-ExSimplex.pdf#page=6}{\textbf{\textcolor{blue}{Basic Variable}}}: \emph{In the Simplex method, a variable that is expressed explicitly in terms of non-basic variables in a dictionary.}%
\item \href{https://www.math.purdue.edu/~yipn/421/2024-08-27-ExSimplex.pdf#page=6}{\textbf{\textcolor{blue}{Non-basic Variable}}}: \emph{In the Simplex method, a variable that is set to zero in a given iteration of the algorithm.}%
\item \href{https://www.math.purdue.edu/~yipn/421/2024-08-27-ExSimplex.pdf#page=9}{\textbf{\textcolor{blue}{Pivot Operation}}}: \emph{The process of exchanging a basic and a non-basic variable in the Simplex method, updating the dictionary to move to a new vertex.}%
\item \href{https://www.math.purdue.edu/~yipn/421/2024-08-27-ExSimplex.pdf#page=26}{\textbf{\textcolor{blue}{Dictionary}}}: \emph{A representation of a linear programming problem where basic variables are expressed as functions of non-basic variables. It is updated during the Simplex method iterations.}%
\item \href{https://www.math.purdue.edu/~yipn/421/2024-08-27-ExSimplex.pdf#page=32}{\textbf{\textcolor{blue}{Auxiliary Problem}}}: \emph{A modified linear programming problem introduced to find an initial feasible solution when the origin is not feasible in the original problem.}%
\end{itemize}

%
\end{itemize}%
\newpage

%
\section{Problem Types}%
\label{sec:ProblemTypes}%
\begin{itemize}[leftmargin=*]%
\subsection{linear programming fundamentals}%
\label{subsec:linearprogrammingfundamentals}%
\begin{itemize}[leftmargin=*]%
\item \href{https://www.math.purdue.edu/~yipn/421/2024-08-27-ExSimplex.pdf#page=32}{\textbf{\textcolor{blue}{Finding an Initial Feasible Solution}}}: \emph{The Simplex method requires an initial feasible solution. If the origin is not feasible, an auxiliary problem is introduced to find a feasible solution to the original problem.}%
\item \href{https://www.math.purdue.edu/~yipn/421/2024-08-27-ExSimplex.pdf#page=26}{\textbf{\textcolor{blue}{Using Dictionaries to Represent Solutions}}}: \emph{The Simplex method uses dictionaries to represent the current solution and to perform pivot operations efficiently. Each dictionary corresponds to a vertex of the feasible region.}%
\end{itemize}

%
\subsection{advanced linear programming}%
\label{subsec:advancedlinearprogramming}%
\begin{itemize}[leftmargin=*]%
\item \href{https://www.math.purdue.edu/~yipn/421/2024-08-27-ExSimplex.pdf#page=32}{\textbf{\textcolor{blue}{Finding an Initial Feasible Solution}}}: \emph{The Simplex method requires an initial feasible solution. If the origin is not feasible, an auxiliary problem is introduced to find a feasible solution to the original problem.}%
\item \href{https://www.math.purdue.edu/~yipn/421/2024-08-27-ExSimplex.pdf#page=30}{\textbf{\textcolor{blue}{Handling Infeasible Origins}}}: \emph{The origin $(0, 0)$ may not satisfy all constraints in a linear programming problem. Techniques like introducing an auxiliary problem are used to find a feasible starting point.}%
\item \href{https://www.math.purdue.edu/~yipn/421/2024-08-27-ExSimplex.pdf#page=17}{\textbf{\textcolor{blue}{Optimizing in Higher Dimensions}}}: \emph{The Simplex method is extended to solve linear programming problems with more than two variables, requiring iterative pivoting operations to move between vertices of the feasible region in higher-dimensional space.}%
\item \href{https://www.math.purdue.edu/~yipn/421/2024-08-27-ExSimplex.pdf#page=26}{\textbf{\textcolor{blue}{Using Dictionaries to Represent Solutions}}}: \emph{The Simplex method uses dictionaries to represent the current solution and to perform pivot operations efficiently. Each dictionary corresponds to a vertex of the feasible region.}%
\end{itemize}

%
\end{itemize}%
\newpage

%
\section{Algorithm Solutions}%
\label{sec:AlgorithmSolutions}%
\begin{itemize}[leftmargin=*]%
\subsection{linear optimization techniques}%
\label{subsec:linearoptimizationtechniques}%
\begin{itemize}[leftmargin=*]%
\item \href{https://www.math.purdue.edu/~yipn/421/2024-08-27-ExSimplex.pdf#page=6}{\textbf{\textcolor{blue}{Simplex Method}}}: \emph{Iteratively move from one vertex of the feasible region to another by setting non-basic variables to zero and choosing a new set of (non)basic variables to improve the objective function value. This corresponds to moving from vertex to vertex in the feasible region.}%
\item \href{https://www.math.purdue.edu/~yipn/421/2024-08-27-ExSimplex.pdf#page=32}{\textbf{\textcolor{blue}{Auxiliary Problem}}}: \emph{When the origin is not feasible, introduce a new variable $x_0$ and modify the constraints to create an auxiliary problem that maximizes $-x_0$. Solve this auxiliary problem using the Simplex method. If the optimal solution has $x_0 = 0$, then the solution is feasible for the original problem.}%
\end{itemize}

%
\end{itemize}%
\newpage

%
\end{document}